\documentclass[11pt]{article}
\usepackage[margin=1in]{geometry}
\usepackage{amsmath,amssymb,amsthm}
\usepackage[T1]{fontenc}
\usepackage{lmodern}
\usepackage[hidelinks]{hyperref}

\newcommand{\F}{\mathbb F}
\newcommand{\Z}{\mathbb Z}
\newcommand{\R}{\mathbb R}
\newcommand{\C}{\mathbb C}

\title{Literature Status for arXiv:1207.0975, Question 4.1}
\author{fa\_banach\_001, agent\_lane\_16}
\date{2026-07-03}

\begin{document}
\maketitle

\paragraph{Status.}
This is a literature-already-answered packet.  The answer is known in a later
paper and is not claimed as a new proof from this run.

\paragraph{Source question.}
Fritz, Netzer, and Thom ask in arXiv:1207.0975, Question 4.1, whether for
\[
  \Gamma=\F_2\times\F_2
\]
the universal group-C*-norm function
\[
  \Z\Gamma\ni a\longmapsto \|a\|_u\in\R
\]
is computable.

\paragraph{Supporting answer.}
Goldbring and Sinclair prove in arXiv:2509.15086, Corollary
\texttt{FNTquestion}, that for every \(n\geq 2\), including \(n=\infty\), the
standard presentation of \(C^*(\F_n\times \F_n)\) is not computable.  They then
state explicitly that the case \(n=2\) answers the question posed by Fritz,
Netzer, and Thom.

\paragraph{Identification.}
For the standard generators of \(\F_2\times\F_2\), generated points in the
standard presentation of \(C^*(\F_2\times\F_2)\) are precisely finite rational
\(*\)-polynomials, equivalently rational group-ring elements, in those
generators.  A computable standard presentation would therefore give an
algorithm that approximates \(\|a\|_u\) for each such \(a\) to arbitrary
precision.  Conversely, the source question's computability of
\(a\mapsto\|a\|_u\) is exactly the norm-computability part of that standard
presentation.  Thus Goldbring--Sinclair's corollary gives a negative answer to
Question 4.1.

\paragraph{Scope.}
This packet concerns the \(F_2\times F_2\) computability question in Section 4
of arXiv:1207.0975.  It does not address the separate question in the same
paper about \(SL_3(\Z)\).

\begin{thebibliography}{9}
\bibitem{FNT}
T. Fritz, T. Netzer, and A. Thom,
\emph{Can you compute the operator norm?},
Proc. Amer. Math. Soc. 142 (2014), 4265--4276; arXiv:1207.0975.

\bibitem{GS}
I. Goldbring and T. Sinclair,
\emph{On definability of C*-tensor norms},
arXiv:2509.15086.
\end{thebibliography}

\end{document}
